\chapter{Bản đồ màn hình}

Mọi việc của thầy cô diễn ra trên một màn hình ba cột: thanh bên trái, khung chat ở giữa và khung đề bên
phải. Khung đề chỉ mở khi thầy cô đang xem một đề.

\anh[0.95\linewidth]{T/04-ban-do}{
  \diem{1}{0.175}{0.075}
  \diem{2}{0.175}{0.2}
  \diem{3}{0.175}{0.37}
  \diem{4}{0.175}{0.62}
  \diem{5}{0.175}{0.92}
  \diem{6}{0.205}{0.5}
  \diem{7}{0.45}{0.88}
  \diem{8}{0.70}{0.06}
  \diem{9}{0.71}{0.93}
}

\begin{bangtra}{L{1cm}L{4cm}Y}{\dau{Số} & \dau{Chỗ} & \dau{Dùng để}}
  \sobuoc{1} & \nut{Đoạn chat mới} & Bắt đầu một việc mới với Kriky. \\
  \sobuoc{2} & Bốn mục chính & \nhan{Bảng theo dõi}, \nhan{Danh sách lớp học}, \nhan{Các bài kiểm tra},
  \nhan{Ngân hàng câu hỏi}. \\
  \sobuoc{3} & \nhan{ĐOẠN CHAT} & Các đoạn chat cũ, xếp theo ngày. Bấm để mở lại. \\
  \sobuoc{4} & \nhan{TÀI LIỆU} & Sách, tài liệu thầy cô đã tải lên (Chương~\ref{chuong:tai-lieu}). \\
  \sobuoc{5} & Tên thầy cô & Menu tài khoản, đăng xuất. Dưới tên là tên bản Kriky. \\
  \sobuoc{6} & Khung chat & Lời nhắn của thầy cô và câu trả lời của Kriky. \\
  \sobuoc{7} & Ô nhắn & Gõ yêu cầu cho Kriky rồi bấm \nut{Gửi}. \\
  \sobuoc{8} & Khung đề & Đề đang xem: các câu, lời giải, nút sửa. \\
  \sobuoc{9} & \nut{Duyệt đề} & Duyệt đề khi đã ưng (Chương~\ref{chuong:duyet}). \\
\end{bangtra}

\machnho{Kéo đường biên giữa hai cột để nới rộng cột nào thầy cô đang đọc.}
